\indent “ಅವನು ಊಟ ಮಾಡಿದನು” ಎಂಬ ವಾಕ್ಯದಲ್ಲಿ “ಮಾಡಿದನು” ಎಂಬ ಶಬ್ದವನ್ನು ನೋಡೋಣ. ಇದು ಭೂತಕಾಲದ ಕ್ರಿಯೆಯನ್ನು ಗೊತ್ತುಪಡಿಸುವ ಕ್ರಿಯಾಪದ.\par

\indent ಮಾಡಿದನು = ಮಾಡು + ದ + ಅನು; ಇಲ್ಲಿ “ಮಾಡು” ಎಂಬ ಧಾತುವಿಗೆ “ಅನು” ಎಂಬ ಆಖ್ಯಾತ ಪ್ರತ್ಯಯ ಸೇರಿ, ಮಾಡಿದನು ಆಗಿರುವುದು. ಆದರೆ, ಧಾತುವಿಗೂ ಆಖ್ಯಾತ ಪ್ರತ್ಯಯಕ್ಕೂ ನಡುವೆ ಬಂದಿರುವ ‘ದ’ ಎಂಬ ಪ್ರತ್ಯಯವು ಕಾಲ ಸೂಚಕ ಪ್ರತ್ಯಯವಾಗಿರುವುದು. ಇದನ್ನು ಹಿಂದಿನ ಪ್ರಕರಣದಲ್ಲಿ ತಿಳಿದಿದ್ದೇವೆ. “ಮಾಡು” ಧಾತುವಿನ ಮತ್ತೊಂದು ರೂಪ “ಮಾಡಿದ” ಎಂಬ ಶಬ್ದ.\par

\indent ಮಾಡಿದ = ಮಾಡು + ದ + ಅ. ಇಲ್ಲಿ ಅಂತ್ಯದಲ್ಲಿರುವ ಪ್ರತ್ಯಯ “ಅ” ಮತ್ತು, ನಡುವೆ ಬಂದಿರುವ ಕಾಲ ಸೂಚಕ ಪ್ರತ್ಯಯ “ದ”. ಇದೇ ರೀತಿಯ ಹಲವು ಉದಾಹರಣೆಗಳನ್ನು ನೋಡೋಣ—\par

\begin{kbvderiv}
\vspace{4pt plus 3pt minus 1pt}
\begin{center}\begin{adjustbox}{max width=\textwidth}
{\renewcommand{\arraystretch}{1.22}\setlength{\tabcolsep}{4pt}
\begin{tabular}{l c l}
({\kbvlatin i}) ನೋಡಿದ & {\color{kbvgrey}$=$} & ನೋಡು + ದ + ಅ \\
({\kbvlatin ii}) ಹಾಡಿದ & {\color{kbvgrey}$=$} & ಹಾಡು + ದ + ಅ \\
\end{tabular}}
\end{adjustbox}\end{center}
\vspace{4pt plus 3pt minus 1pt}
\end{kbvderiv}

\indent (ಇಲ್ಲಿ “ದ” ಎಂಬುದು ಭೂತಕಾಲ ಸೂಚಕ ಪ್ರತ್ಯಯ)\par

\begin{kbvderiv}
\vspace{4pt plus 3pt minus 1pt}
\begin{center}\begin{adjustbox}{max width=\textwidth}
{\renewcommand{\arraystretch}{1.22}\setlength{\tabcolsep}{4pt}
\begin{tabular}{l c l}
({\kbvlatin iii}) ಓಡುವ & {\color{kbvgrey}$=$} & ಓಡು + ವ + ಅ \\
({\kbvlatin iv}) ತಿನ್ನುವ & {\color{kbvgrey}$=$} & ತಿನ್ನು + ವ + ಅ \\
\end{tabular}}
\end{adjustbox}\end{center}
\vspace{4pt plus 3pt minus 1pt}
\end{kbvderiv}

\indent (ಇಲ್ಲಿ “ವ” ಎಂಬುದು ವರ್ತಮಾನ ಕಾಲ ಸೂಚಕ ಪ್ರತ್ಯಯ)\par

\begin{kbvderiv}
\vspace{4pt plus 3pt minus 1pt}
\begin{center}\begin{adjustbox}{max width=\textwidth}
{\renewcommand{\arraystretch}{1.22}\setlength{\tabcolsep}{4pt}
\begin{tabular}{l c l}
({\kbvlatin v}) ನಡೆಯುವ & {\color{kbvgrey}$=$} & ನಡೆ + ಉವ + ಅ \\
({\kbvlatin vi}) ಕುಣಿಯುವ & {\color{kbvgrey}$=$} & ಕುಣಿ + ಉವ + ಅ \\
\end{tabular}}
\end{adjustbox}\end{center}
\vspace{4pt plus 3pt minus 1pt}
\end{kbvderiv}

\indent (ಇಲ್ಲಿ “ಉವ” ಎಂಬುದು ಭವಿಷ್ಯತ್ಕಾಲ ಸೂಚಕ ಪ್ರತ್ಯಯ) ಮೇಲಿನ ಆರು ಉದಾಹರಣೆಗಳಲ್ಲೂ ಪ್ರತಿಯೊಂದು ಶಬ್ದದ ಅಂತ್ಯದಲ್ಲಿ “ಅ” ಎಂಬ ಪ್ರತ್ಯಯವಿರುವುದು. ಈ ಪ್ರತ್ಯಯವನ್ನು ಕೃತ್‌ ಪ್ರತ್ಯಯ ಎನ್ನುತ್ತೇವೆ. ಹೀಗೆ ಕೃತ್‌ ಪ್ರತ್ಯಯವನ್ನು ಅಂತ್ಯದಲ್ಲಿ ಉಳ್ಳದ್ದೇ ಕೃದಂತ.\par

\noindent\hspace*{2.2em}{[}ಕೃತ್‌ + ಅಂತ = ಕೃದಂತ{]}.\par
